\documentclass[11pt]{article}
\usepackage{fontspec}
\usepackage[french]{babel}
\usepackage{amsmath}
\usepackage{amssymb}
\usepackage{amsthm}
\usepackage{mathrsfs}
\usepackage{tikz-cd}
\usepackage{enumitem}
\usetikzlibrary{babel}

\newcommand{\printedpage}[1]{\marginpar{\footnotesize p.~#1}}
\newcommand{\msp}[1]{\marginpar{\footnotesize [#1]}}
\newcommand{\Ch}{\underline{\mathrm{Ch}}}
\newcommand{\Ty}{\underline{\mathrm{Ty}}}
\newcommand{\dashto}[1]{\mathrel{\overset{#1}{\rule[0.6ex]{2em}{0.4pt}}}}

\begin{document}

\printedpage{89}
% head: A. Grothendieck : Dérivateurs, Ch. VII, Édition des Universités de Montpellier II et Paris VII
% foot: Version 22 décembre 2001

\msp{page 143}
Dans $1^o$ et $3^o$, on veut prouver que certaines flèches ($p$ et $\bar{f}$) sont à la fois \textit{propre} et \textit{lisse} - ce qui est vraiment idéal ! Je dis qu'il suffit à chaque fois de prouver l'un des deux, que l'autre s'en déduit (comme `cas particulier') par dualité. Pour ceci, notons
\[
\tau \longmapsto \tau^{-}
\]
l'automorphisme (\textit{pas anti}-automorphisme) de $\Ty$ qui consiste à changer de signe toutes les flèches de transition - c'est donc, au niveau des objets, l'unique automorphisme du monoïde libre Ty, qui échange $\Delta_+$ et $\Delta_-$. Au système projectif $\underline{t}$ dans $\Ty$ correspond le système projectif $\underline{t}^-$. Ceci dit, on a un isomorphisme de catégories
\[
\boxed{\Ch^{\underline{t}}_{\infty}(X^o) \ \simeq\ (\Ch^{\underline{t}^-}_{\infty}(X))^o} \tag*{$(^{35})$} % Reading resolved against the supplied PDF; see corrections.csv.
\]
\footnotetext[35]{cf. page 140 pour le passage de $\underline{t}$ à $\underline{t}^o$.}
et le diagramme commutatif
\[
\begin{tikzcd}
\Ch^{\underline{t}}_{\infty}(X^o) \arrow[r, "\sim"] \arrow[d, "p^{\underline{t}}_{X^o}"] & \Ch^{\underline{t}^-}_{\infty}(X)^o \arrow[d, "(p^{\underline{t}^-}_{X})^o"] \\
X^o \times X^o \arrow[r, equal] & (X \times X)^o.
\end{tikzcd}
\]
Ainsi, dire que $p^{\underline{t}}_{X}$ \texttt{[plutôt} $p^{\underline{t}^-}_{X}$\texttt{]} est lisse, i.e. que $(p^{\underline{t}}_{X})^o$ \texttt{[plutôt} $(p^{\underline{t}^-}_{X})^o$\texttt{]} est propre, équivaut à dire que $p^{\underline{t}}_{X^o}$ est propre.

\msp{page 144}
Il nous reste donc, à ce sujet, que deux faits à vérifier : la lissité de $p : \Ch_{\infty}(X) \longrightarrow X \times X$, et celle de $\bar{f} : \overline{X'} \longrightarrow X$.

a) Lissité de $p : \Ch_{\infty}(X) \longrightarrow X \times X$. On se donne une flèche
\[
w : (a,b) \longrightarrow (a',b'), \qquad \text{i.e.} \quad \left\{ \begin{array}{lccc} u : & a & \longrightarrow & a' \\[1ex] v : & b & \longrightarrow & b', \end{array} \right.
\]
sur la base, un objet $c'$ au dessus du but, i.e.
\[
c' \in \Ch_{\infty}(X; a', b'),
\]
et on regarde la catégorie des flèches dans $\Ch_{\infty}(X)$, de but $c'$, qui relèvent $w$,
\[
\begin{tikzcd}
\underset{\text{\scriptsize variable}}{c} \arrow[r, "\varphi"] \arrow[d, dash] & \underset{\text{\scriptsize fixe}}{c'} \arrow[d, dash] \\
(a,b) \arrow[r, "w"] & (a',b'),
\end{tikzcd}
\]
\printedpage{90}
% head: A. Grothendieck : Dérivateurs, Ch. VII, Édition des Universités de Montpellier II et Paris VII
% foot: Version 22 décembre 2001
i.e. la catégorie des couples $(c,\varphi)$, où
\[
\left\{ \begin{array}{l}
c \in \Ch_{\infty}(X; a, b) \\[1ex]
\varphi : c \longrightarrow c' \ \ (\in \Ch_{\infty}(X; a', b')) \\[1ex]
\text{une flèche induisant } \begin{array}[t]{l} u \ \text{ sur les sources} \\ v \ \text{ sur les buts,} \end{array}
\end{array} \right.
\]
\[
\begin{tikzcd}
a \arrow[rrrr, dotted, dash, "c"] \arrow[d] & {} \arrow[d, dotted] & {} \arrow[d, dotted] & {} \arrow[d, dotted] & b \arrow[d] \\
a' \arrow[rrrr, dash, "c'"] & {} & {} & {} & b'.
\end{tikzcd}
\]

\msp{page 145}
\begin{flushright}
{\footnotesize 9.11.90}
\end{flushright}

Pour la question précédente, relative aux flèches dans le diagramme
\[
\begin{tikzcd}
& X \arrow[d, "i"] & \\
& \Ch_{\infty}(X) \arrow[d, "p"] \arrow[ddl, "\sigma"'] \arrow[ddr, "\beta"] & \\
& X \times X \arrow[dl, "\mathrm{pr}_1"] \arrow[dr, "\mathrm{pr}_2"'] & \\
X & & X,
\end{tikzcd}
\]
il faut préciser absolument la déscription de $\Ch_{\infty}(X)$ \textit{dans} Cat, donc il faut préciser le système projectif $\underline{t} = (t_n)_{n \in \mathbf{N}}$. Nous allons prendre en effet % [sic] déscription -> description, misspelling
\[
t_n = (\underbrace{\Delta_+ \Delta_-}_{\tau_0})^{2n},
\]
avec les morphismes de transition décrits par ailleurs
\[
\underbrace{t_{n+p}}_{=\tau_0^{2n+2p}} \ \longrightarrow \ \underbrace{t_n}_{\tau_0^{n}}, % [sic] \tau_0^{n} -> \tau_0^{2n}, since t_n = (\Delta_+\Delta_-)^{2n} = \tau_0^{2n}
\]
\printedpage{91}
% head: A. Grothendieck : Dérivateurs, Ch. VII, Édition des Universités de Montpellier II et Paris VII
% foot: Version 22 décembre 2001
étant donnés par $\delta_p * \mathrm{id}_{\tau_o^{2n}} * \delta_p$, où $\delta_p : \tau_0^p \longrightarrow 1$ est la dégénérescence de $\tau_0^p$. Ainsi on peut considérer $\Ch_{\infty}(X)$ comme étant la catégorie des chemins biïnfinis dans $X$ % [sic] \tau_o^{2n} -> \tau_0^{2n}, letter o printed in the index instead of 0
\[
\cdots \longrightarrow a_{-5} \longleftarrow \overbrace{a_{-4} \longrightarrow a_{-3} \xleftarrow{u_{-2}} \overbrace{a_{-2} \xrightarrow{u_{-1}} a_{-1} \xleftarrow{u_0} \overbrace{a_0}^{\Ch_0} \xrightarrow{u_1} a_1 \xleftarrow{u_2} a_2}^{\Ch_1} \xrightarrow{u_3} a_3 \longleftarrow a_4}^{\Ch_2} \longrightarrow a_5 \longleftarrow \cdots,
\]
avec les $u_i$ des identités pour $|i|$ assez grand.

La sous-catégorie (pleine)
\[
\Ch^{W}_{t_n}(X) \ = \ \Ch^{W}_{(\Delta_+\Delta_-)^{2n}}
\]
s'identifie à celle pour laquelle on a
\begin{quote}
$u_i$ une identité pour $i \notin [-2i+1, 2i]$, i.e. les seuls $u_i$ `significatifs' sont les $4i$ flèches $u_{-2i+1}$, $u_{-2i+2}$, $\ldots$, $u_{2i}$. % [sic] i (in [-2i+1,2i], 4i, u_{-2i+1}, u_{-2i+2}, u_{2i}) -> n, the index of t_n is n
\end{quote}

Mais il sera plus commode de regarder $\Ch_{\infty}(X)$ comme la limite inductive d'une
\msp{page 146}
famille à deux indices de sous-catégories pleines
\[
\Ch_{n,m}(X), \qquad n, m \in \mathbf{Z},
\]
en définissant cette catégorie comme \textit{vide} si $n > m$, et pour $n \leq m$ comme formée des objets
\[
\left\{ \begin{array}{l}
(a_i, u_i)_{i \in \mathbf{Z}} \in \Ch^{W}_{n,m}(X) \iff u_i \text{ une identité pour } i \notin [n, m-1] \\[1ex]
\overset{\mathrm{id}}{\cdots}\, a_{n-1} \dashto{\mathrm{id}} \underbrace{a_n \dashto{u_n} a_{n+1} \cdots \dashto{} a_{m-1} \dashto{u_{m-1}} a_m}_{\text{\scriptsize `partie utile' du chemin}} \dashto{\mathrm{id}} a_{m+1} \overset{\mathrm{id}}{\cdots}\,.
\end{array} \right.
\]
Donc pour $n = m$, on trouve la sous-catégorie pleine isomorphe à $X$, formée des chemins infinis `totalement dégénérés', i.e. tels que \textit{toutes} les $u_i$ soient des identités. (Cette sous-catégorie de $\Ch_{\infty}$ ne dépend donc pas du choix de $n = m$.)
\[
\xrightarrow{\mathrm{id}} \underbrace{a_{-1}}_{=\, a} \xleftarrow{\mathrm{id}} \underbrace{a_0}_{=\, a} \xrightarrow{\mathrm{id}} \underbrace{a_1}_{=\, a} \xleftarrow{\mathrm{id}}
\]
Tous les sommets d'un tel chemin sont égaux à un même objet de $X$, et les flèches de transition entre eux sont identiques.

La structure de $\Ch^{W}_{n,m}(X) \ \simeq \ \Ch^{W}_{\tau_{n,m}}(X)$ ne dépend que de la longueur $m-n$ de $\tau_{n,m}$ et de la \textit{parité} de $n$. Si $n$ est pair, la direction initiale de $\tau_{n,m}$ est positive, et
\[
\tau_{n,m} = \underbrace{\Delta_+ \Delta_- \cdots}_{m-n \text{ facteurs en alternance}} = \left\{ \begin{array}{ll}
(\Delta_+ \Delta_-)^{\frac{m-n}{2}} & \text{si } m \text{ pair} \\[1ex]
(\Delta_+ \Delta_-)^{\frac{m-n}{2}-1} \Delta_+ & \text{si } m \text{ impair}
\end{array} \right.
\]
\texttt{[plutôt} $\frac{m-n-1}{2}$ \texttt{au lieu de} $\frac{m-n}{2} - 1$\texttt{]}, tandis que pour $n$ impair c'est juste l'inverse
\printedpage{92}
% head: A. Grothendieck : Dérivateurs, Ch. VII, Édition des Universités de Montpellier II et Paris VII
% foot: Version 22 décembre 2001
\msp{page 147}
\[
\tau_{n,m} = \underbrace{\Delta_- \Delta_+ \cdots}_{m-n \text{ facteurs alternés}} = \left\{ \begin{array}{ll}
(\Delta_- \Delta_+)^{\frac{m-n}{2}} & \text{si } m \text{ impair} \\[1ex]
(\Delta_- \Delta_+)^{\frac{m-n}{2}-1} \Delta_- & \text{si } m \text{ pair}
\end{array} \right.
\]
\texttt{[plutôt} $\frac{m-n-1}{2}$ \texttt{au lieu de} $\frac{m-n}{2} - 1$\texttt{]}. La direction initiale est conforme à la parité de $n$, i.e. égale à $(-1)^n$, tandis que la direction terminale est opposée à la parité de $m$, i.e. égale à $(-1)^m$ \texttt{[plutôt} $(-1)^{m+1}$\texttt{]}.

Je note que l'on a $\tau_0 = (\tau_0)^o$, d'où
\[
\underline{t} = \underline{t}^o, \qquad \text{mais } \underline{t} \neq \underline{t}^-,
\]
et pourtant on a des isomorphismes fonctoriels en $X$
\[
\Ch_{\underline{t}}(X) \ \simeq \ \Ch_{\underline{t}^-}(X)
\]
respectant source et but, à cause d'un isomorphisme
\[
\underline{t} \ \simeq \ \underline{t}^- \qquad \text{dans Pro}(\Ty).
\]
On peut dire p. ex. que le \textit{foncteur translation} $T^p$ d'amplitude $p$ (pour $p \in \mathbf{Z}$) est défini sur $\Ch^{W}_{\underline{t}}(X)$, par
\[
\left\{ \begin{array}{l}
T^p(((a_i),(u_i))_{i \in \mathbf{Z}}) = ((b_j),(v_j))_{j \in \mathbf{Z}} \\[1ex]
\text{avec} \left\{ \begin{array}{lcl} b_j & = & a_{j-p} \\[1ex] v_j & = & u_{j-p} \end{array} \right.
\end{array} \right.
\]
et que
\[
\begin{aligned}
T^p &: \left\{ \begin{array}{lcl} \Ch_{t}(X) & \xrightarrow{\sim} & \Ch_{-t}(X) \\[1ex] \Ch_{-t}(X) & \xrightarrow{\sim} & \Ch_{t}(X) \end{array} \right\} \ \text{si } p \text{ impair} \\
T^p &: \left\{ \begin{array}{lcl} \Ch_{t}(X) & \xrightarrow{\sim} & \Ch_{t}(X) \\[1ex] \Ch_{-t}(X) & \xrightarrow{\sim} & \Ch_{-t}(X) \end{array} \right\} \ \text{si } p \text{ pair.}
\end{aligned} % Reading resolved against the supplied PDF; see corrections.csv.
\]

\msp{page 148}
En d'autres termes, on peut formellement regarder la catégorie somme
\[
\Ch_{\pm}(X) = \Ch_{+}(X) \amalg \Ch_{-}(X)
\]
et l'opération translation $T$ comme un automorphisme qui échange les deux morceaux, et $T^p$ n'est autre que son itéré $p$-ème, qui échange les morceaux $+$ et $-$ si $p$ impair, qui les stabilise si $p$ pair.

\printedpage{93}
% head: A. Grothendieck : Dérivateurs, Ch. VII, Édition des Universités de Montpellier II et Paris VII
% foot: Version 22 décembre 2001

Il y a aussi l'automorphisme involutif de passage au chemin opposé, qui échange source et but, et que je définis (conformément à l'isomorphisme canonique identique
\[
\underline{t} \xrightarrow{\sim} \underline{t}^o \quad )
\]
comme la symétrie $S$
\[
\left\{ \begin{array}{l}
S((a_i),(u_i)) = ((b_j),(u_j)) \ \ \text{avec} \\[1ex]
\begin{array}{lcl} b_j & = & a_{-j} \\[1ex] u_j & = & u_{-j+1}. \end{array}
\end{array} \right.
\]
Si je le regarde comme opérant sur $\Ch^{W}_{\pm}(X)$, il est lié au foncteur $T$ par
\[
\left\{ \begin{array}{rcl}
S^2 & = & \mathrm{id} \\[1ex]
STS^{-1} & = & T^{-1},
\end{array} \right.
\]
de sorte que c'est le groupe semidirect de $\mathbf{Z}/2\mathbf{Z}$ par $\mathbf{Z}$ (groupe affine de dimension 1 de l'anneau $\mathbf{Z}$) qui opère sur $\Ch_{\pm}$.

\msp{page 149}
On est embarassé de trouver un automorphisme $T^2$ de la situation $\Ch_{\infty}(X) \longrightarrow X \times X$, sans qu'il soit évident que dans Hot, ce foncteur induise l'identité! Je ne vois aucun moyen de relier un chemin $c$ et $T^2 c$ par une flèche dans $\Ch_{\infty}$, ni même pas un chemin `canonique'. Mais je me dis que cet automorphisme est décrit par un automorphisme, non pas du système projectif % [sic] embarassé -> embarrassé, misspelling
\[
(\tau_{n,m})_{n,m \in \mathbf{Z}^o \times \mathbf{Z},\ n \leq m}
\]
(où $\mathbf{Z}^o$ désigne l'ensemble ordonné \textit{opposé} de $\mathbf{Z}$), mais du \textit{pro-objet} de $\Ty$ qu'il définit (qui est bien sûr `le même' pour tous les systèmes projectifs cofinaux à gauche dans $\Ty$). Il en résulte déjà que $T^2$ induit l'identité sur l'ind-objet de Cat\,hot, $\underline{\underline{\mathrm{Ch}}}_{\infty}(X)$, et a fortiori sur l'ind-objet correspondant de Hot, ou de $\mathrm{Hot}^{\wedge}_{0}$, ou enfin pour l'objet de $\mathrm{Hot}^{\wedge}_{0}$ lui-même. De même pour les $\Ch_{\infty}(X; a, b)$. Il s'ensuit, revenant à $\Ch_{\infty}(X)$ lui-même et aux $\Ch_{\infty}(X; a, b)$, que l'automorphisme induit par $T^2$ est l'identité dans $\mathrm{Hot}^{\wedge}_{0}$. J'espère que $\mathrm{Hot} \longrightarrow \mathrm{Hot}^{\wedge}_{0}$ est
\msp{page 150}
tout au moins fidèle, de sorte que ça prouvera que $T^2$ induit l'identité dans Hot également.

Je reviens aux questions $1^o$ à $4^o$ d'il y a quelques jours.

Pour $1^o$), il faut prouver que
\[
p : \Ch_{\infty}(X) \longrightarrow X \times X
\]
est lisse et propre et une Hot-fibration. Par dualité, $p$ propre résultera de $p$ lisse. Mais $\Ch_{\infty}(X)$ est $\varinjlim$ des $\Ch_{n,m}(X)$, avec $n \longrightarrow -\infty$, $m \longrightarrow +\infty$, et où on peut se borner à $n$ et $m$ \textit{pairs.} Or on voit que pour $n$, $m$ pairs, i.e. $\tau_{n,m}$ à flèche initiale directe, à flèche finale

\end{document}
